\section{Related Work}
\label{sec:related_work}

\paragraph{LLM-based multi-agent collaboration.}
Large language models have recently been used as autonomous agents that collaborate through natural language~\citep{han2024llmmultiagent}. 
Early multi-agent collaboration frameworks~\citep{li2023camel, metagpt, chatdev} focus on leveraging large language models in diverse roles and interaction patterns to solve tasks.
Beyond task-oriented settings, a parallel line of work explores LLM-based agents as simulacra of human behavior, where agents represent individuals~\citep{park2023generative,Wang2024AgentVerse,hua2023warandpeace, yang2024oasis}. 
More recently, agent collaboration has expanded toward agent-to-agent interaction, where agents operate as independent entities.

\paragraph{Privacy constraints in LLMs.}
Prior work on privacy in large language models has investigated how LLMs can inadvertently expose sensitive information and how this risk can be mitigated. 
Studies have shown that pretrained LLMs can infer personal attributes or leak training data during inference, highlighting privacy vulnerabilities inherent in model capabilities~\citep{staab2023beyond}. 
Privacy surveys in LLMs systematically categorize these threats and review mitigation strategies such as differential privacy, data sanitization, and secure inference mechanisms~\citep{YAO2024100211, miranda2024preservingPrivacySurvey}. 
Differential privacy has also been applied to LLM prompt learning and in-context learning to provide formal privacy guarantees while balancing utility~\citep{duan2023flocks, tang2024privacy}. 

\paragraph{Privacy constraints in multi-agent collaboration.}
As LLM-based agents increasingly operate as autonomous entities, recent work has examined security and privacy issues in agent-to-agent (A2A) interactions~\citep{a2a_protocol_what_is_a2a}. 
More recent work shows that LLM-based multi-agent systems introduce new vulnerabilities, including message interception, manipulation, and leakage of sensitive contextual information during agent-to-agent conversations~\citep{gomaa2025converse, he2025red}. 
Several frameworks further propose architectural safeguards, such as monitoring or sentinel agents, to enforce security and policy constraints during agent interactions~\citep{gosmar2025sentinel, nakamura2025terrarium}. 
However, this line of work largely treats privacy and security as properties of the communication infrastructure, rather than examining how privacy constraints reshape collaborative reasoning and coordination. 
